% !PRINTED 305
\weberpage{305}

Die Aufgabe der Auffindung aller dieser metacyklischen
Gleichungen ist also darauf zur\"uckgef\"uhrt, alle metacyklischen
Theiler der homogenen Congruenzgruppe $S$ zu finden.

Nehmen wir als Beispiel den Fall $p = 3$, $k = 2$, fragen
also nach den metacyklischen Gleichungen 9\textsuperscript{ten} Grades, so handelt
es sich nach diesem Satze um die Gruppe $S$ der nach dem
Modul 3 genommenen linearen Substitutionen

\weq{12}{\sigma = \begin{pmatrix} a, & b \\ c, & d \end{pmatrix},}

mit denen wir uns, unter etwas anderem Gesichtspunkte, im \S.\,66
besch\"aftigt haben.

Wir haben dort zun\"achst einen Theiler $E$ von $S$ betrachtet,
der aus allen Substitutionen $\sigma$ mit der Bedingung $a\,d - b\,c$
$\equiv 1$ (mod 3) besteht, und haben ausserdem zwei Substitutionen

\weq{13}{\begin{pmatrix} a, & b \\ c, & d \end{pmatrix}, \quad \begin{pmatrix} -a, & -b \\ -c, & -d \end{pmatrix}}

zu einem einzigen Elemente zusammengefasst. Die so specialisirte
Gruppe war vom Grade 12 und hatte einen Normaltheiler
vom Index 3:

\[ G = \begin{pmatrix} 1, & 0 \\ 0, & 1 \end{pmatrix}, \begin{pmatrix} 0, & 1 \\ -1, & 0 \end{pmatrix}, \begin{pmatrix} -1, & 1 \\ 1, & 1 \end{pmatrix}, \begin{pmatrix} 1, & 1 \\ 1, & -1 \end{pmatrix}, \]

aus dem $E$ so zusammengesetzt war:

\[ E = G + \begin{pmatrix} 1, & 0 \\ 1, & 1 \end{pmatrix} G + \begin{pmatrix} 1, & 0 \\ -1, & 1 \end{pmatrix} G. \]

In der Gruppe $S$ gelten aber die beiden Substitutionen (13)
als verschieden, und ausserdem m\"ussen noch die Substitutionen
$\sigma$, deren Determinante $a\,d - b\,c \equiv -1$ (mod 3) ist, dazu
genommen werden, wodurch sich der Grad der Gruppe auf 48
erh\"oht. $G$ erweitert sich durch die Aenderung der Vorzeichen
aller Elemente zu einer Gruppe 8\textsuperscript{ten} Grades und $E$ zu einer
Gruppe 24\textsuperscript{sten} Grades, von der das erweiterte $G$ Normaltheiler
ist. Hier ist nun die ganze Gruppe $S$ \gesperrt{metacyklisch}, wie man
aus folgender Zusammensetzung sieht, in der die erweiterte
Gruppe $G$ mit $S_2$ und die erweiterte Gruppe $E$ mit $S_1$ bezeichnet
 ist:

\[ S_4 = \begin{pmatrix} 1, & 0 \\ 0, & 1 \end{pmatrix}, \begin{pmatrix} -1, & 0 \\ 0, & -1 \end{pmatrix} \]
%CONT>>
